%=======================================================================
%  Paper V — Section 9 (Conclusions)   FINAL
%
%  Requires the macros listed at the head of Section 2.
%
%  Superseded values against ver10:
%      "2e-14 per cent"          -> 2.1e-14 in relative terms
%      0.11 R centroid offset    -> 0.1125 R
%      26.68 % -> 26.6 %, the value the run records (26.59 %)
%      I = 16.889 %              -> 16.9 %   (grid floor, Section 6.1)
%      C = 1.507                 -> 1.51 +/- 0.09
%      Bennu I = 0.000 %         -> < 10^-3 % (negative I observed)
%      Dx_COM = 16.5 m           -> 16.5 +/- 2.9 m; Lowry is 1.75 sigma
%      +/- 0.435 %               -> +/- 0.434 %
%      s = -1.14                 -> -1.141 +/- 0.018, eleven planes
%      correlation 0.95          -> 0.95 with n = 4, p ~ 0.03
%      rho_rest = 1855 +/- 8     -> 1858 +/- 27  (Eq. budget)
%      Delta M = 2.75e9, 7.7 %   -> (2.75 +/- 0.48)e9, 7.7 +/- 1.3 %
%      sigma to six decimals     -> five (Section 6.1)
%      turning point at 0.25 km  -> REMOVED; argument rebuilt on the
%                                   11.2 sigma compositional exclusion
%      Arrokoth lower bound 290  -> central ~235, range 155-600
%      7.5 kg/m3 per metre       -> 10.0 at the reference plane, and
%                                   0.28 for rho_rest
%      "no diagnostic detects a misplaced plane"
%                                -> three of four respond; A fails
%      2.2e-14 per cent          -> relative, and a check on the
%                                   quadrature only
%      "19 to 59 times"          -> add the working-resolution figure
%      "confirmed and conservative"
%                                -> confirms the ORDERING; stability and
%                                   accuracy are different quantities
%
%  Added, from results absent in ver10:
%      conclusion 5  exclusion of the mass dipole at 21 sigma
%      conclusion 4  decomposition of the stability ratio; S as the
%                    signal-to-noise of the excess mass
%      conclusion 9  reporting rule keyed to the larger piece
%      closing       the first Richardson condition does not transfer
%
%  All values confirmed. Resolved here: the Arrokoth 395/410 against
%  390/419.7 discrepancy (buried facets retained in the resolution series);
%  the Itokawa-specific envelope, 1.35 %, and the ratio 12.5 against it; the
%  235 kg/m3 run and the three conflicting published density constraints;
%  standard errors on the ellipsoid ratios; the completed residual table and
%  the omega = 0 control. Only the Zenodo DOI is still a placeholder.
%=======================================================================

\section{Conclusions}
\label{sec:conclusions}

This paper takes the interior density inversion of \citet{richardson2014}
and \citet{kanamaru2019} as an object of study rather than as a tool, and
asks what it determines. Five bodies were examined---Itokawa, Eros,
Bennu, 67P/Churyumov--Gerasimenko and Arrokoth---together with eight
synthetic contact binaries constructed so that the strength of the signal
could be varied while everything else was held fixed, and a homogeneous
ellipsoid whose surface gravity has a closed form. The main conclusions
are as follows.

\begin{enumerate}

%---------------------------------------------------------------
\item \emph{Regions must be built by exact plane clipping, and the
natural alternative fails silently.} Closing a region by fanning its
boundary edges to an interior point yields a watertight mesh whose volume
matches the sum of its own tetrahedra to $2.1\times10^{-14}$ in relative
terms, and which encloses a radial cone rather than the plane-cut lobe
intended. On a unit sphere cut at $x = 0.5$ the two solids differ by a
factor $1.600$ in volume and by $0.1125\,R$ in centroid; on Itokawa the
error reversed the sign of the recovered contrast. A volume-closure test
cannot detect it, because it compares the mesh against the same
tetrahedra that built it. Five layers are required, of which three can
fail in ways the others cannot: a predicate test, requiring every vertex
of the region to satisfy the inequality defining it, which the cone fails
outright; a per-component manifold test, which catches a surface torn
along the cut that every volume-based test passes; and an audit of the
identity $\Dmass = (\densbulk-\densrest)\Vtot = \Ddens\VR$ on every
solution, which operates on the output rather than on the geometry and
would have caught the cone error there, the two volumes differing by the
same factor $1.600$. The clipped construction reproduces closed-form cap
and slab volumes to $0.09\%$, limited by the polyhedral approximation of
the sphere itself, and reproduces the independently published lobe volume
of 67P, $26.6\%$ against $27\%$ \citep{jorda2016}.

%---------------------------------------------------------------
\item \emph{The density recovered by the inversion is resolution-stable
over the meshes tested, and the residual sensitivity enters through the
plane detector rather than through the objective function.} This does not
follow from the stability of the dispersion established in Paper~IV,
since the minimum of a stable function need not itself be stable, and was
therefore tested separately. Three meshes of Itokawa sharing a dividing
plane, spanning a $64\%$ increase in facet count, return $\denshead$
unchanged; because no sub-grid refinement is applied, that bounds the
variation at half a scan step, $\SI{12.5}{\kilo\gram\per\cubic\metre}$ or
$0.45\%$, rather than resolving it. Three meshes of Arrokoth at a fixed
plane, spanning $173\%$, likewise return an identical
$\SI{395}{\kilo\gram\per\cubic\metre}$; that series retains the buried
facets of Section~\ref{sec:methods:shapes}, since it compares meshes with
one another rather than reporting an absolute value, which is why its
densities are $395$ and $410$ rather than the $390$ and $419.7$ of
Table~\ref{tab:arrokothdensity}. The informative case is the
coarsest Itokawa mesh, on which the detector moves the plane by
$\SI{9.4}{\metre}$; at the local gradient that displaces $\denshead$ by
about $3.4\%$, comparable to its own curvature interval, and $\densrest$
by $0.14\%$, where it adds nothing to the budget. Resolution sensitivity
therefore matters for the quantity the method does not constrain and not
for the one it does, and a study that fixed the plane externally would
not see it at all. Three meshes are not a convergence sequence, and no
claim of convergence in the analytical sense is made.
% Resolved: the resolution series retains the buried facets and the table
% excludes them; both files now say so (Sections 4.7 and 9).

%---------------------------------------------------------------
\item \emph{Of five bodies, only Itokawa yields a dispersion improvement
above the mesh-resolution sensitivity envelope of Paper~IV.} At the
detected neck the improvement is $16.9\%$ with
$\denshead/\densrest = 1.51 \pm 0.09$, a factor $2.9$ above the upper
edge of the $0.9$--$5.9\%$ envelope, which measures discretization
sensitivity rather than statistical noise. That upper edge is attained on
a different body. Against the value the envelope takes on Itokawa itself,
$1.35\%$ over the ladder from $2{,}000$ to $42{,}000$ facets with quadric
decimation, the improvement is $12.5$ times larger, and that is the
meaningful comparison.
Eros, Bennu, 67P and Arrokoth return $0.418\%$, below $10^{-3}\%$,
$0.574\%$ and $1.272\%$, all within or below that envelope and all
returning contrasts consistent with unity once the curvature interval is
propagated through the mass constraint---Eros, the nearest to a
detection, at $1.018 \pm 0.026$, or $0.7$ intervals from unity. Each is
explicable: Eros is measured to be nearly uniform, Bennu's reported
heterogeneity is radial and invisible to a plane cut, 67P's
centre-of-mass position is stated by \citet{jorda2016} to be inexplicable
by two uniform lobes, and Arrokoth has no measured mass. At the dividing
plane $x = \SI{0.150}{\kilo\metre}$ used by \citet{kanamaru2019} we
recover a contrast of $1.467$ against their $1.460$; the agreement tests
implementations rather than physics, since both minimize a functional of
the surface geopotential, and our own interval on that contrast is
$\pm0.09$, two orders of magnitude wider than the difference. The same
plane gives a centre-of-mass offset of $16.5 \pm \SI{2.9}{\metre}$,
consistent with the $\SI{21+-12}{\metre}$ of \citet{lowry2014}.

%---------------------------------------------------------------
\item \emph{Where the method responds strongly, the quantity it
constrains is the background density of the body---and equivalently the
mass anomaly---not the density of the region; and that holds in the
strong-signal regime rather than of the method as such.} The mass
constraint makes $\Dmass = \Ddens\VR$ equal to $\Mtot-\densrest\Vtot$
exactly, so a statement that one is invariant is the same statement about
the other; these are one result and not two, and the pair
$(\densR,\VR)$ remains degenerate along $\Ddens\VR = \mathrm{const}$.
Across the thirteen-plane Itokawa sweep the head density varies by
$\pm38.3\%$ while the background density varies by $\pm0.434\%$, and in
coefficients of variation the ratio is $82$. Most of that is not a
result. Equation~\eqref{eq:excessamplify} fixes a factor of $11.8$ of it as
algebra, present whatever the objective does, and the subtraction of a
nearly constant background offsets a further $2.4$, which is also algebra.
What remains, the factor of $16.9$ by which $\Dmass$ is steadier than
$\Ddens$, measures compensation between contrast and region volume; over
the eleven planes the composition permits the figures are $32.1$, $11.78$,
$2.77$ and $7.52$, and over the nine below the softer bound $27.3$, $11.88$,
$2.92$ and $6.71$. It is the compensation that is the result. On 67P and
Arrokoth the
order of stability reverses, and the controlled synthetic experiment
resolves the discrepancy through the exponent
$s = \mathrm{d}\log\lvert\Ddens\rvert/\mathrm{d}\log\VR$, which takes the
value $-1$ when the background density is held fixed. Bodies with a weak
response show a density difference that changes sign, so no exponent
exists; those with a moderate response give $s \simeq -2$; only the
strongest gives $-1.22$, against Itokawa's $-1.141 \pm 0.018$ over the
eleven planes that span the same factor $1.907$ in region volume as the
synthetic family. The correlation between signal strength and exponent is
$0.95$, but over the four synthetic bodies that have an exponent, three of
them clustered in signal strength, so the direction of the effect is
supported and its functional form is not determined. Because the
synthetic bodies are uniform by construction, this establishes a property
of the objective function rather than a physical inference. For Itokawa
the constrained quantities are
$\densrest = 1858 \pm \SI{27}{\kilo\gram\per\cubic\metre}$ and
$\Dmass = (2.75 \pm 0.48)\times10^{9}$~kg, or $7.7 \pm 1.3\%$ of the
total, the interval in both cases being dominated by the curvature
sensitivity interval and not by the plane-placement scatter, which is
three times smaller and is not an uncertainty. Equivalently, the
displacement statistic $\Sstat$ of Equation~\eqref{eq:sstat} equals
$\Dmass/\delta(\Dmass)$ identically, so its value of $5.8$ on Itokawa is
the statement that the excess mass is determined to $17\%$.

%---------------------------------------------------------------
\item \emph{The conserved quantity is the excess mass and not the mass
dipole, and the two are separated by twenty-one standard errors.} The
alternative hypothesis---that the objective holds fixed
$\Dipole = \Dmass\,(x_R-x_{\mathrm{cf}})$, which is the quantity the YORP
determination of \citet{lowry2014} constrains---makes a different and
quantitatively separated prediction, because the geometric factor is not
constant along a sweep. Over the eleven admissible planes
$\mathrm{d}\log(x_R-x_{\mathrm{cf}})/\mathrm{d}\log\VR = -0.244$, fixed by
the shape model alone, so dipole conservation requires $s = -0.756$; the
measured $-1.141 \pm 0.018$ lies on the opposite side of $-1$. Stated
directly, $\mathrm{d}\log\Dipole/\mathrm{d}\log\VR = -0.385 \pm 0.018$,
twenty-one standard errors from the zero that conservation requires, and
the conclusion does not depend on which subset of planes is used: the same
calculation gives $10.9\sigma$ over all thirteen and $16.8\sigma$ over the
nine. This matters beyond the internal question, because it makes the
agreement with \citet{lowry2014} a genuinely independent check rather than
two readings of one quantity.

%---------------------------------------------------------------
\item \emph{Within the two-region planar model, dispersion minimization
does not constrain its own dividing plane, and the demonstration requires
no extrapolation.} The dispersion at the minimum falls monotonically as
the plane advances, from $0.07549$ to $0.06373$ across the thirteen
Itokawa planes, an improvement rising from $12.9\%$ to $26.4\%$ without a
single reversal. At the outermost plane the recovered head density is
$\SI{5000}{\kilo\gram\per\cubic\metre}$, which exceeds the LL-chondrite
grain density of $3540 \pm \SI{130}{\kilo\gram\per\cubic\metre}$---a
zero-porosity ceiling---by $11.2$ standard deviations, and the empirical
block-density bound of $3220 \pm \SI{220}{\kilo\gram\per\cubic\metre}$ by
$8.1$. Since the improvement is strictly increasing and the admissible set
is an interval bounded above, the improvement attains its maximum over
that set at a boundary fixed by the meteorite collection rather than by
the minimization, and no statement about the behaviour of $\potvarn$
beyond the sweep is required. Converted through the measured gradient the
two bounds give $\xsplit \le \SI{0.2069+-0.0047}{\kilo\metre}$ and
$\le \SI{0.1935+-0.0110}{\kilo\metre}$. The reason follows from
conclusion~4: with $\densrest$ held nearly fixed,
$\densR = \densrest + \Dmass/\VR$ diverges as $\VR \to 0$, and nothing in
the objective opposes the divergence. Cubic and quartic fits place a
turning point near $\SI{0.254}{\kilo\metre}$, but we do not rest anything
on it: the figure is an extrapolation beyond the last plane computed, the
measured gradient extrapolates instead to $0.262$--$\SI{0.278}{\kilo
\metre}$, the position coincides with the plane at which the scan ceiling
begins to truncate the recovered contrast, and it is not established that
an interior turning point exists at all, the objective having a finite
limit as the region shrinks to a point mass at the tip.

%---------------------------------------------------------------
\item \emph{The plane must come from outside the method, and the two
external criteria available are not equally strong.} The neck detector
returns $\xsplit = \SI{160.8}{\metre}$ from the cross-sectional profile
alone, and the material bound excludes planes beyond
$\SI{194+-11}{\metre}$; the detected neck falls $\SI{33}{\metre}$ inside
the bound, about three times the bound's own uncertainty. That the two do
not contradict one another is reassurance and not corroboration, because
the compositional criterion is one-sided: it excludes planes that are too
far out and says nothing about planes that are too far in, and every plane
below $\SI{194}{\metre}$ satisfies it, which is $58\%$ of the range swept.
The objective's own preference among admissible planes is real but not one
the method can defend: it prefers the block-density boundary to the
detected neck by $4.6$ percentage points of improvement, which lies inside
the mesh-resolution envelope. The saddle prominence, which should be
published so that a reader can judge the provenance of the plane,
separates accepted from refused bodies by more than two orders of
magnitude but carries no information about the strength of the result:
Arrokoth has the highest prominence in the study, $0.490$, and returns
$1.272\%$, while Itokawa has the third, $0.170$, and returns $16.9\%$.

%---------------------------------------------------------------
\item \emph{On a body without a measured mass, the assumed bulk density
can reverse the sign of the inferred contrast, not merely change its
scale.} For Arrokoth the contrast runs from $0.9999$ at an assumed
$\SI{250}{\kilo\gram\per\cubic\metre}$ to $0.9159$ at $500$, and at fixed
assumed mass the plane sweep carries the excess mass from $+2.5\%$ to
$-11.6\%$ of the total, passing through zero at a region fraction of
$61.2\%$---five percentage points from the plane the detector selects. The
method cannot say which lobe is denser. A fourth run at the central published estimate of \citet{keane2022},
$\SI{235}{\kilo\gram\per\cubic\metre}$, returns $1.066$: the sense of the
answer reverses between $235$ and $250$, inside the published one-sigma
range of $155$ to $600$ and below the bound of $290$ that
\citet{spencer2020} derive from the requirement that mutual gravity
exceed the centrifugal acceleration. What the inversion reports about this
body is close to a restatement of what was assumed about it. The converse suggests a diagnostic: on Itokawa a fourfold change
in assumed mass moves the contrast by $11\%$ and never near unity. The
outcome must be read as the position of the resulting interval relative to
unity rather than as its width, since on the raw change Itokawa is the
more sensitive of the two bodies tested and the separation appears only
per e-fold of assumed mass. We put this forward as a candidate rather than
an established test: it rests on two bodies that differ in several
respects besides signal strength, and it is of use precisely where the
mass is unknown and there is nothing to calibrate it against.
% Resolved: 235 with a 1-sigma range 155-600 is Keane et al. (2022); the
% 290 lower bound is Spencer et al. (2020) under the no-tension condition;
% McKinnon et al. (2020) give 250-500. All three are now cited in Sections
% 2.9 and 4.7, and the 235 run has been performed.

%---------------------------------------------------------------
\item \emph{Seven failure modes produce plausible-looking numbers, and no
single diagnostic separates a result from an artefact.} A region too small
flattens the minimum until the density interval reaches $22\%$ of its own
value; a complement too small sharpens it, and the tightest curvature
bound in this study, $\pm\SI{15}{\kilo\gram\per\cubic\metre}$, belongs to
the one run known to be wrong---an artefact of the coordinate scanned,
since propagating it through the mass constraint at a region fraction of
$88.4\%$ gives $\pm\SI{115}{\kilo\gram\per\cubic\metre}$, or $3.83\%$, the
worst determination anywhere in the study. A minimum at a scan edge
reported an apparent $13.6\%$ improvement from a value pinned to the scan
floor. Facets lying inside a shape model delivered as two unjoined shells,
$0.943\%$ of the surface area, produced a crossing of unity on Arrokoth
that is not there. And a plane misplaced onto a flank of Eros returned a
contrast of $0.876 \pm 0.039$, excluding unity by more than three
intervals, with every geometric check passing. Four numerical diagnostics
were tested against that last case and they do not agree: the
amplification factor fails outright, ranking the artefact as the strongest
signal in the study at $\lvert\ampfac\rvert = 9.12$ against Itokawa's
$12.0$, because $\ampfac$ depends on which piece is called the region; the
displacement statistic ranks it correctly at $\Sstat = 2.86$ against
$5.8$, and is the only diagnostic invariant under that relabelling, but it
places the artefact above all four genuine nulls; the envelope comparison
rejects it by a margin of $6\%$; and the decisive argument is the
provenance of the plane, which is not available from the output. A
validity window on region volume fraction was considered and rejected, the
degradation having no threshold in the data. What replaces it is
procedural: place the plane at a true saddle, publish its prominence,
propagate the curvature interval through the mass constraint before
reading it, report the sensitivity of the reported quantity to moving the
plane, and test facet centroids for containment whenever a model has more
than one component.

%---------------------------------------------------------------
\item \emph{The quantity to report is the density of the larger of the two
pieces, together with the excess mass and the amplification factor.}
Differentiating the mass constraint gives
$\lvert\delta\densrest\rvert/\lvert\delta\densR\rvert = f/(1-f)$, so the
density of the larger piece is always the better determined and the
threshold is exactly $f = \tfrac{1}{2}$. On Itokawa, Eros and 67P the
region is the smaller piece and $\densrest$ is the constrained quantity;
on Arrokoth the detector places $66.3\%$ of the volume in the region, so
there $\densrest$ is the density of the smaller piece and is worse
determined by a factor of $1.83$. Which piece is called the region is a
labelling choice and carries no physics: it changes neither uncertainty,
only which density a reporting rule names. Report $\Dmass$ alongside,
since it is the same statement and is the quantity that survives the
relabelling, and report $\ampfac$ with the plane at which it was
evaluated, since a fractional range in $\densrest$ corresponds to one in
$\Dmass$ that is \emph{larger} by that factor---on Itokawa, $2.95\%$
against $0.251\%$---and since a large $\ampfac$ is a warning rather than a
precision. The sensitivity to plane placement must likewise be quoted for
the quantity being reported: at the Itokawa reference plane it is
$\SI{10.0}{\kilo\gram\per\cubic\metre\per\metre}$ for $\denshead$, rising
to $55$ at the outer planes, against
$\SI{0.28}{\kilo\gram\per\cubic\metre\per\metre}$ for $\densrest$ and
$0.18\%$ per metre for $\Dmass$.

%---------------------------------------------------------------
\item \emph{Measured against an analytic ellipsoid, the geopotential
dispersion is more accurate than the mean surface slope at every facet count
tested, by a factor near $8$ at the $49{,}152$-facet resolution at which the
Itokawa and Eros results are computed, and by $19\pm4$ to $59\pm18$ over the
coarser range $2{,}000$ to $32{,}000$ facets.} The working-resolution figure
is the one a summary should quote, and it carries two qualifications: it is a
single shape and spin state rather than the average of nine, and the
dispersion error against which it is formed, $-0.02\%$, is quoted to one
significant figure, so the ratio is determined only to within roughly $6$ to
$12$. This test needs no assumption about which mesh is the
truth: the surface gravity of a homogeneous ellipsoid has a closed form,
and our evaluation of it reproduces the sphere limit to
$2.2\times10^{-14}$ in relative terms---a check on the numerical
quadrature, not on the polyhedral code, which cannot reproduce a sphere to
any such precision. Two qualifications belong with the factor. The
advantage narrows systematically with resolution, the two errors
converging as $N^{-1.26}$ and $N^{-0.91}$ so that their ratio falls as
$N^{-0.35}$; extrapolated to the $49{,}152$-facet meshes on which the
Itokawa and Eros results are computed it is of order $16$ to $20$, and
that is the figure a summary should quote. And the comparison is with
Paper~IV's \emph{ordering} of the two quantities, which it confirms
independently of the assumption Paper~IV could not test; it is not a
confirmation that Paper~IV's factor of $3$ to $14$ was conservative, since
mesh-to-mesh variation measures stability and the ellipsoid measures
accuracy, and a quantity can be stable and biased. The same experiment
shows vertex clustering and quadric edge collapse to be closely comparable
in accuracy, with mean absolute slope errors of $3.92\%$ and $4.05\%$, so
the numbers reported in Papers~III and~IV, which used clustering, are not
in question; quadric collapse is adopted here for topology rather than
accuracy. The ellipsoid also bounds its own use: decimating over the same
range moves the dispersion by about $0.18$ percentage points on these
smooth bodies against the $0.9$--$5.9\%$ Paper~IV measured on real ones,
so surface roughness dominates the envelope and the analytic test
validates its ordering rather than its magnitude.
% Resolved: standard errors are now given in Tables 8.3 and 8.4, and the
% decimator comparison is made on the paired differences, +0.019 +/- 0.012
% per cent on the dispersion and +0.12 +/- 0.20 on the slope, neither
% significant. The 49,152-facet row was abandoned as disproportionately
% expensive; one configuration completed and gives a ratio near 8, quoted
% in Section 8.3 as an indication.

\end{enumerate}

Three qualifications frame these results.

The relaxation premise on which the whole approach rests is unverified for
every body examined, and on Itokawa, where three of the four conditions of
\citet{richardson2014} are supported, the first fails: they tabulate its
mean rotational potential as $2.4\%$ of the mean gravitational potential
and their global bulk-density inference on it returns
$\SI{330}{\kilo\gram\per\cubic\metre}$ against a measured $1950$. The
single strong signal reported here is thus on the body where the premise
is weakest. Two things bear on how much weight that should carry.
Section~\ref{sec:discussion:conditioning} shows that the failing condition
is the condition under which a \emph{bulk-density} inversion is not
identically flat, and that it has no counterpart in a fixed-mass
two-region inversion: substituting the mass constraint into the
geopotential makes the surface potential affine in the single free
parameter, $\Upot_i = \alpha_i + \beta_i\densR$, and $\alpha$ and $\beta$
are proportional only if $\Mtot = 0$, so an interior stationary point
survives even in the limit of no rotation. The mass constraint fixes the
scale that the rotational term fixes in their problem. That is a statement
about conditioning and not about physics: it removes a reason to expect
failure and does not establish that the surface has relaxed. Empirically,
the inference is insensitive to the assumed mass over a factor of four and
agrees with a YORP determination that makes no appeal to relaxation at
all---though that determination is itself only $1.75$ standard deviations
from zero, and its interval contains the centre-of-mass offset at every
plane of the sweep, so it corroborates the sign and order of magnitude of
the anomaly and discriminates neither the plane nor the value. Against
this, Itokawa's residual dispersion after the best fit, $7.2\%$, is large
in absolute terms and larger than Eros's $3.4\%$; that comparison is not
controlled for shape, and the residuals for the other four are $98.7\%$ on Arrokoth,
$99.4\%$ on 67P, $99.6\%$ on Eros and $100.0\%$ on Bennu, so Itokawa's
$83.1\%$ stands apart from a tightly grouped set that includes two other
bilobate bodies. The $\omega = 0$ control of Section~\ref{sec:discussion}
has also been run: with the rotational term removed the scan still returns
an interior minimum, at $\SI{2950}{\kilo\gram\per\cubic\metre}$ with an
improvement of $17.2\%$, so the conditioning argument is confirmed in the
output; the contrast moves by $7.1\%$, which is larger than the error
budget and should be read beside it.
% Both done: Table 8.5 is complete and the omega = 0 control is reported
% (v10_s8_tests.json).

The model admits two uniform regions divided by a plane, a form known to
be insufficient for two of the bodies studied and shown by
\citet{kanamaru2019pss} to become largely unconstrained if relaxed to
three regions on Itokawa. A mass anomaly does not specify the region that
carries it, which is why the head--neck question raised by
\citet{scheeres2008} and revisited by \citet{lowry2014} and
\citet{kanamaru2019pss} remains open in general: structures placing
comparable excess mass at comparable distance from the centre of figure
produce comparable centre-of-mass offsets and comparable dispersion
reductions, and no surface measurement resolves them. The head and neck
models of this paper are not an instance of that degeneracy and should not
be cited as one---they place excess masses differing by a factor of $2.2$
and are separated by the objective by a factor of $7.3$ in improvement and
by $5.8$ against $2.0$ on $\Sstat$, so on the global surface the head
model is preferred and the neck model does not clear the resolution
envelope at all. What the two models do show is that the invariance
established here holds under a change of plane position within a fixed
region shape and not under a change of region shape, the two returning
background densities $\SI{85}{\kilo\gram\per\cubic\metre}$ apart.

The recovered densities are quantized at the
$\SI{25}{\kilo\gram\per\cubic\metre}$ scan step, with no sub-grid
refinement, and the consequences are visible throughout: two planes
$\SI{0.8}{\metre}$ apart returning the same $\denshead$, two non-monotone
entries in the centre-of-mass column, and an improvement that can come out
marginally negative where the grid does not contain the uniform solution,
which is an arithmetic impossibility for the underlying problem. Since
$\potvarn^{2}$ is a ratio of quadratics in $\densR$ once the mass
constraint is substituted, its stationarity condition is linear with a
single root, and solving it in closed form would remove the grid from the
error budget entirely and yield the curvature interval as the roots of a
quadratic rather than as a fitted quantity. We have not done so, and it is
the clearest improvement available to a future implementation.

The natural continuation is a calibrated detection limit. The four
weak-signal cases reported here are individually explicable, but the study
cannot say whether one strong signal in five reflects the rarity of large
interior contrasts or the sensitivity floor of the method. Such a test
cannot be performed by planting a contrast in a body and leaving its shape
alone, since the objective is built from the shape and never reads a field
computed beforehand; the premise requires that the surface have relaxed
onto an equipotential of the true interior, so the shape must be built
before the inversion is applied. Our attempt to do so did not converge
(Section~\ref{sec:discussion:rate}), and its mode of failure is
instructive: the plane holding a fixed volume fraction moved as the
surface relaxed, so the shape and the dividing surface co-evolve, and a
fixed point requires solving for both together---which is the coupling
Section~\ref{sec:plane} reports from the other direction. A stable
relaxation scheme is a piece of work in its own right. It would settle
both the detection limit and the interpretation of the synthetic
experiment, which as performed here varies shape rather than density, and
we regard it as a paper rather than an appendix.

\vspace{1em}
\noindent\textbf{Data availability.} All code, driver scripts and derived
tables, including the five-layer self-test and the mass-balance audit over
all thirteen Itokawa planes, are archived at Zenodo
(\href{https://doi.org/10.5281/zenodo.23003040}{doi:10.5281/zenodo.23003040})
under CC~BY~4.0, and the archived configuration records the scan interval
used at every plane so that the widening described in
Section~\ref{sec:itokawa:sweep} can be audited rather than taken on trust.
Development is tracked at
\url{https://github.com/praponkan/shape-based-surface-gravity}. Shape
models are not redistributed and are identified by their published
references in Section~\ref{sec:methods}; Table~\ref{tab:shapes} gives the
facet counts, volumes, masses and spin periods actually used, and the
deposit records the archive that holds each model.
% DOI reserved 2026-09-28: 10.5281/zenodo.23003040, filled here and in
% Section 2.9. Confirm before submission that the published record is the
% one that produced the numbers.
